\chapter{Carbohydrate-Deficient Transferrin (CDT)}

\section*{What Is This Test?}
Carbohydrate-deficient transferrin measures transferrin molecules with reduced carbohydrate content. Sustained heavy alcohol exposure can increase these transferrin forms, making CDT a biomarker of recent chronic alcohol use rather than an immediate alcohol level.

\section*{Alternative Names}
CDT; Percent CDT; Transferrin Isoforms; Alcohol Biomarker.

\section*{Why This Test Is Ordered}
CDT may support assessment of sustained alcohol exposure, monitoring in selected treatment or safety programs, and evaluation when the history is uncertain. It should be interpreted with clinical assessment and other biomarkers such as gamma-glutamyl transferase or carbohydrate-deficient transferrin-related testing.

\section*{How This Test Is Performed}
A blood sample is analyzed by an electrophoretic or chromatographic method that separates transferrin isoforms. Results are usually reported as a percentage of total transferrin.

\section*{Preparation, Risks, and Limitations}
No fasting is usually required. Venipuncture is the usual risk. Liver disease, genetic transferrin variants, pregnancy, iron deficiency, and some other conditions can affect results. CDT is less useful for occasional drinking, binge drinking, or a single recent exposure.

\section*{Understanding Results}
An elevated percentage supports sustained heavy alcohol exposure but is not proof of alcohol use or alcohol-use disorder. A normal result does not exclude harmful drinking. Interpretation should include timing, liver tests, history, and laboratory-specific cutoffs.

\section*{References}
World Health Organization guidance on alcohol biomarkers and clinical assessment of alcohol use.

\clearpage
\section*{Regulatory Status and Authorization Date}
Regulatory status applies to the specific assay, instrument, manufacturer, and intended use rather than to the test name alone. No single approval date can be assigned to this test category. For a commercial assay, verify the current product in the relevant regulator's database: the U.S. Food and Drug Administration (FDA) clearance, approval, or authorization; India's Central Drugs Standard Control Organisation (CDSCO) licence or permission; the European Union In Vitro Diagnostic Regulation (EU IVDR) CE certificate issued through the applicable conformity-assessment route; or the United Kingdom Medicines and Healthcare products Regulatory Agency (MHRA) registration and UKCA/CE status. Laboratory-developed tests may not have an individual FDA, CDSCO, EU IVDR, or MHRA approval date.
\clearpage
